% 第四章待用数学素材；本文件不接入 main.tex，不是第四章完整正文。
% 从第三章原3.5.5—3.5.6迁出，保留原数学内容。
% 正式整合第四章时，应按该章结构改写衔接，并更新 eq:ch3-19/20 标签。
% 第三章提供 M_i、d_i、Q_i、P_base、W、a_r、D_r 与共享频率域定义。

\subsection{部分配置的能量下界}

求解过程中，若集合\(J\)中的Stage已经确定模式，而其余Stage只保留非空候选集\(\widehat M_i\subseteq M_i\)，则固定部分按其实际活动能量计入，其余部分仍取独立最小值，得到该分支的乐观下界

\begin{equation}
L(J)=P_{\rm base}W+\sum_{i\in J}Q_i(m_i)d_i(m_i)
+\sum_{i\notin J}\min_{m\in\widehat M_i}Q_i(m)d_i(m).
\label{eq:ch3-19}
\end{equation}

这样，完整配置可以直接计算能量，尚未确定全部模式的分支则由\(L(J)\)估计最低可能代价。安全见证提供可实现的能量上界，可靠上下界闭合时便可证明最优。实际计算需区分独立可靠界与求解器报告值，其处理沿用第三章求解状态的定义。这里的独立能量比较用于构造下界，而不是直接删除其他模式，因为这些模式在时长、设备占用或频率兼容性上仍可能不同。

\subsection{固定冲突方向的无热判定}

固定配置后，剩余选择包括资源使用顺序和各Stage的起点。同一执行单元上的Stage对，以及使用不同CPU单元但属于同一频率域的异频Stage对，都必须确定先后关系；同频的不同CPU单元Stage仍可重叠。将全部互斥对的先后方向与原有依赖边合并，得到有向图\(H\)。任何合法安排都对应一组这样的方向，满足完整方向图的安排也就满足了资源与频率域互斥。

判断这一固定方向能否满足期限，可从最早执行安排入手。若\(H\)存在有向环，沿环相加的先后约束会要求零不小于一组严格正时长之和，故该方向不可行。若\(H\)无环，则按拓扑序依次安排Stage：其最早起点取请求到达时刻与所有前驱完成时刻中的最大值，即

\begin{equation}
\eta_i=\max\left\{a_{r(i)},\max_{j\in\operatorname{Pred}_H(i)}[\eta_j+d_j(m_j)]\right\}.
\label{eq:ch3-20}
\end{equation}

无前驱时，式\eqref{eq:ch3-20}只取请求到达时刻。按拓扑序归纳可知，任意满足该方向图的起点均有\(s_i\ge\eta_i\)。因此，这一方向无热可行当且仅当图无环且全部\(\eta_i+d_i(m_i)\le D_{r(i)}\)：最早安排若已超期，继续等待也无法满足期限；若均按期完成，最早安排本身就是无热可行安排。网格版本在每一步将最大值向上对齐到\(\Delta\)格点，再以真实时长传播到后继，使每个Stage的起点都满足依赖与网格要求。

温度要求使这一判定还需向前一步。最早安排虽然逐个Stage采用最早起点，却未必在每个时刻产生最低注热；推迟Stage会在部分时刻减少功率，又在后续时刻增加功率。因此，最早安排通过热检查时可作为安全见证，出现超温时则仍需考虑同一方向下的其他等待时序。求解器需要寻找低能量配置的完整安全安排，并对无法实现的选择给出相应排除依据。
